\subsection{Mapping based on Hamming distance}
\label{sec:mapping}

\paragraph{Presentation}MiMyCS is a software that performs reference-based lossless compression of next-gen sequencing datasets~\cite{DeMoor2024}.
It relies on mapping each read against a known reference and storing the mapping position along with potential substitution errors.
The mapping stage is implemented with PiM acceleration by partitioning the reference across memory banks and sending reads to map to only a few well-chosen processing units.
We re-implement this mapping procedure using the BPL library.

\paragraph{Limitations encountered} We process the reference and need to store the resulting indexing information for all future mapping executions.
However, the BPL does not support such feature at the moment.
Therefore, we had to design an additional class to manage persistent storage and implement it for all architectures supported by the BPL.
For UPMEM, the persistent indexing tables are allocated in the MRAM of each DPU.
As for Multicore, we allocate all tables required, store them in a vector, and adds an index argument to the task so that each instance can retrieve its assigned tables.

\paragraph{PiM specifics} To perform the Hamming distance computations as efficiently as possible on the UPMEM architecture, we write manually a few lines of assembly instructions.
This part of the task thus uses target-specific code.

\paragraph{Benchmarks} We consider a real whole genome sequencing dataset obtained with the Illumina technology, known as accession code SRR14724532 on the European Nucleotide Archive portal.
We map this dataset against the human genome GRCh38.
For PiM, we use 8 ranks as this is the minimum required to fit the entire reference genome across all memory banks.
We report the execution time of the two stages: indexing the reference, and mapping all the reads of the dataset.

\begin{table}[htbp]
    \centering
    \caption{Results of the mapping benchmark}
    \label{tab:mapping-results}
    \begin{tabular}{lrrrr}
    \toprule
     &\multicolumn{2}{c}{Multicore time (s)} & \multicolumn{2}{c}{PiM time (s)} \\
    \cmidrule(l{3pt}r{3pt}){2-3}\cmidrule(l{3pt}r{3pt}){4-5}
    & Standard & BPL & SDK & BPL \\
    \midrule
    Indexing reference & 33.1 & 49.4 & 20.4 & 34.2 \\
    Mapping reads & 104.6 & 142.2 & 59.1 & 92.8 \\
    \bottomrule
    \end{tabular}
\end{table}

\paragraph{Results} We present results in \tablename~\ref{tab:mapping-results}.
For both the indexing and the mapping, we see that the BPL introduces a non-negligible overhead and is 63\% slower compared to the optimized SDK version.
The overhead is for the most part due to the serialization process when retrieving the task results back to the host.
Indeed, this involves several memory copies: first on the DPUs, to re-aggregate all the tasklets results into a single buffer, and then on the host, to reformat the transferred buffer into the expected type.
All of these copies do not exist when using the SDK.

\paragraph{Source code} The BPL implementation is publicly available at \url{https://gitlab.inria.fr/pim/org.pim.mimycs-bpl} (GNU AGPLv3).
The original Standard Multicore and SDK implementations are available at \url{https://gitlab.inria.fr/pim/org.pim.srm} (GNU AGPLv3).